% from ega0-18-fr.tex
\oldpage[0\textsubscript{IV-1}]{58}

On dit que la $A$-extension $H$ de $B$ par $L'$ est \emph{déduite de $E$ au
moyen de l'homomorphisme} $w:L\to L'$. Le caractère fonctoriel de la somme
amalgamée en chacun des sommandes montre en outre que si l'on a un morphisme
d'extensions
\[
  \xymatrix{
    0\ar[r] & L\ar[r]\ar[d]^{1_L}\ar@{}[d]|{\wr}
      & E_1\ar[r]\ar[d]^g & B_1\ar[r]\ar[d]^h & 0\\
    0\ar[r] & L\ar[r] & E_2\ar[r] & B_2\ar[r] & 0
  }
\]
on en déduit canoniquement un morphisme d'extensions
\[
  \xymatrix{
    0\ar[r] & L'\ar[r]\ar[d]^{1_{L'}}\ar@{}[d]|{\wr}
      & E_1\oplus_LL'\ar[r]\ar[d]^{g\oplus_L1_{L'}}
      & B_1\ar[r]\ar[d]^h & 0\\
    0\ar[r] & L'\ar[r] & E_2\oplus_LL'\ar[r] & B_2\ar[r] & 0.
  }
\]
En particulier, si $E_1$ et $E_2$ sont des $A$-extensions $A$-équivalentes de
$B$ par $L$, les extensions de $B$ par $L'$ qu'on en déduit au moyen de $w$
sont $A$-équivalentes.

Lorsque l'on a un morphisme (18.2.4.1) de $A$-extensions, il se \emph{factorise
à travers} la $A$-extension $H$ de $B$ par $L'$ déduite de $E$ au moyen de
l'homomorphisme $w:L\to L'$ ($L'$ étant considéré comme $B$-bimodule au moyen
de l'homomorphisme $v:B\to B'$)~: en effet, la définition de la somme amalgamée
montre qu'il existe un $A$-homomorphisme unique $u_0:H\to E'$ de
$A$-bimodules, rendant commutatif le diagramme
\[
  \xymatrix{
    0\ar[r] & L\ar[r]^j\ar[d]^{w_0} & E\ar[r]^f\ar[d]^{j_1}
      & B\ar[r]\ar[d]^{1_B}\ar@{}[d]|{\wr} & 0\\
    0\ar[r] & L'_0\ar[r]^{j_2}\ar[d]^{j_0}\ar@{}[d]|{\wr}
      & H\ar[r]^{f_0}\ar[d]^{u_0} & B\ar[r]\ar[d]^v & 0\\
    0\ar[r] & L'\ar[r]_{j'} & E'\ar[r]_{f'} & B'\ar[r] & 0
  }
\]
avec $u=u_0\circ j_1$, $w=j_0\circ w_0$, $j_1$ et $j_2$ étant les
homomorphismes canoniques~; on vérifie immédiatement que $u_0$ est aussi un
homomorphisme d'anneaux.

Notons enfin les propriétés fonctorielles relatives aux extensions triviales~:
\end{env}

\begin{proposition}[18.2.9]
\label{0.18.2.9-fr}
Soient $B$, $B'$ deux $A$-anneaux, $L$ un $B$-bimodule, $L'$ un
$B'$-bimodule, $v:B\to B'$ un $A$-homomorphisme d'anneaux, $w:L\to L'$ un
homomorphisme de $A$-bimodules tel que $(v,w)$ soit un di-homomorphisme de
bimodules. Alors, il existe un $A$-homomorphisme unique d'anneaux
$u:D_B(L)\to D_{B'}(L')$ rendant commutatif les diagrammes
\[
  \xymatrix@C=2.2pc@R=1.6pc{
    D_B(L)\ar[r]^u & D_{B'}(L') & & D_B(L)\ar[r]^u & D_{B'}(L')\\
    B\ar[u]\ar[r]_v & B'\ar[u] & & L\ar[u]\ar[r]_w & L'\ar[u]
  }
\]
où les flèches verticales sont les injections canoniques.
\end{proposition}

